% !TeX program = lualatex
% Self-contained: no external bibliography, images, or \input files.
\documentclass[11pt,a4paper]{article}
\usepackage[margin=24mm,headheight=14pt]{geometry}
\usepackage{fontspec}
\setmainfont{Latin Modern Roman}
\setsansfont{Latin Modern Sans}
\setmonofont{Latin Modern Mono}
\usepackage{amsmath,amssymb,mathtools}
\usepackage{unicode-math}
\setmathfont{Latin Modern Math}
\usepackage{microtype}
\usepackage{enumitem,xurl,needspace}
\usepackage[dvipsnames]{xcolor}
\usepackage{hyperref}
\hypersetup{unicode=true,colorlinks=true,linkcolor=MidnightBlue,urlcolor=MidnightBlue,
 pdftitle={Problem statement},pdfauthor={Source-annotated compilation},bookmarksnumbered=true}
\usepackage{bookmark,tocloft,titlesec,fancyhdr}
\setlength{\cftsecnumwidth}{4.2em}
\cftsetpnumwidth{2.5em}
\cftsetrmarg{4em}
\titleformat{\section}{\Large\bfseries\raggedright}{\thesection}{0.7em}{}
\pagestyle{fancy}\fancyhf{}
\fancyhead[L]{\small\sffamily Additional UnsolvedMath statements}
\fancyhead[R]{\small Original catalogue IDs}
\fancyfoot[C]{\thepage}
\setlength{\parindent}{0pt}
\setlength{\parskip}{5pt plus 1pt}
\setlength{\emergencystretch}{3em}
\setcounter{tocdepth}{1}\setcounter{secnumdepth}{1}
\setlist[itemize]{leftmargin=3em,itemsep=4pt,topsep=4pt,parsep=0pt}
\allowdisplaybreaks
\newcommand{\N}{\mathbb{N}}
\newcommand{\Z}{\mathbb{Z}}
\newcommand{\Q}{\mathbb{Q}}
\newcommand{\R}{\mathbb{R}}
\newcommand{\C}{\mathbb{C}}
\newcommand{\F}{\mathbb{F}}
\newcommand{\A}{\mathbb{A}}
\newcommand{\PP}{\mathbb{P}}
\newcommand{\E}{\mathbb{E}}
\newcommand{\eps}{\varepsilon}
\newcommand{\dd}{\,\mathrm d}
\newcommand{\sref}[2]{\hyperlink{source:#1:#2}{[S#2]}}
\newif\ifshowcatalogue
\showcataloguetrue
\newcommand{\cataloguescope}[1]{\ifshowcatalogue
 \paragraph{Statement recorded in the supplied source.}
 #1\fi}
\begin{document}

\section{Near-Linear Numbers of Distances from One Point}\label{statement}


\subsection{Definitions and mathematical statement}

\paragraph{Shared notation.}Unless a section says otherwise, $\N=\{1,2,\ldots\}$,
$\Z$ is the ring of integers, $\Q,\R,\C$ are the rational, real, and complex
fields, $[N]=\{1,\ldots,\lfloor N\rfloor\}$, and $\log$ is the natural
logarithm. For real functions of a parameter tending to infinity,
$f\ll g$ and $f=O(g)$ mean $|f|\leq Cg$ eventually for some fixed $C$;
$f\gg g$ reverses this comparison; $f=o(g)$ means $f/g\to0$;
$f\sim g$ means $f/g\to1$; and $f\asymp g$ means both order bounds.
A subscript on $O$ or $\ll$ specifies allowed constant dependence.
The expression $n^{o(1)}$ in an upper bound means growth smaller than
$n^\epsilon$ for each fixed $\epsilon>0$. Local definitions govern any
exception, finite-field convention, or use of the same symbol for another
object. An asymptotic claim is not automatically an effective algorithm.



Distances are Euclidean. Distinct distances mean distinct numerical values; when a distance is said to occur a certain number of times, unordered pairs of distinct points are counted. Congruence allows any Euclidean isometry, including reflection. \sref{statement}{1} \sref{statement}{2}

\paragraph{Statement recorded in the supplied source.}
Given $n$ distinct points $A\subset\mathbb{R}^2$ must there be a point $x\in A$ such that $ \#\{ d(x,y) : y \in A\} \gg n^{1-o(1)}? $ Or even $\gg n/\sqrt{\log n}$? \sref{statement}{1}

\paragraph{Residual task reported by the archive.}

The embedded review (2026-08-17) records: Improve the pinned-distance exponent to n\textasciicircum{}(1-o(1)) or the grid-scale bound. \sref{statement}{1}

\subsection{Short English statement}

Must some point of every large planar configuration see nearly as many different distances as there are points, at the proposed sharp scale?

\subsection{Sources}

{\small

\begin{itemize}

\item[\hypertarget{source:statement:1}{S1}] ULAM.AI, \emph{UnsolvedMath}, supplied \texttt{problems.json}, catalogue entry, statement and background. The embedded literature-review date is 2026-08-17; that is the archive’s review date, not a fresh certification. Dataset landing page: \url{https://huggingface.co/datasets/ulamai/UnsolvedMath}.

\item[\hypertarget{source:statement:2}{S2}] T. F. Bloom, \emph{Erdős Problems}, Problem 604, \url{https://www.erdosproblems.com/604}. Reference from the supplied record; the live problem page was not independently checked for this compilation.

\item[\hypertarget{source:statement:3}{S3}] Erd\H{o}s, Paul, On some problems of elementary and combinatorial geometry. Ann. Mat. Pura Appl. (4) (1975), 99-108. Bibliographic information from the supplied record.

\item[\hypertarget{source:statement:4}{S4}] Erd\H{o}s, Paul, Some of my favourite unsolved problems. Math. Japon. (1997), 527-537. Bibliographic information from the supplied record.

\item[\hypertarget{source:statement:5}{S5}] Katz, Nets Hawk and Tardos, G\'{a}bor, A new entropy inequality for the Erd\H{o}s distance problem. Towards a theory of geometric graphs (2004), 119-126. Bibliographic information from the supplied record.

\end{itemize}

}

\label{end:statement}
\end{document}
